\bookchapter{第六十三章：视图、状态与导航}

\chapterintro{单页面的玩具已经会做了。这一章解决两个“长大”问题：怎么在多个页面之间跳来跳去（导航），以及数据在页面之间怎么流动（状态管理）。学完本章，你就能搭出“列表 详情 编辑”这种经典的多页面 App 骨架。}

\booksection{一、状态管理的“数据流全景”}

SwiftUI 的状态工具不止 @State 一个。先看这张全景图，搞懂“什么时候用哪个”：

\begin{center}
\begin{tabularx}{\linewidth}{|X|X|X|X|}
\hline
\textbf{工具} & \textbf{管什么} & \textbf{谁来用} & \textbf{典型场景} \tabularnewline
\hline
@State & 页面私有的临时状态 & 当前视图 & 计数器、输入框文字 \tabularnewline
@Binding & 借用别人的状态 & 子视图改父视图 & 自定义组件 \tabularnewline
@StateObject / @Observable & 跨页面共享的逻辑 & 视图模型（ViewModel） & 数据加载、业务逻辑 \tabularnewline
@EnvironmentObject / .environment & 全局共享的状态 & 整个 App & 用户登录态、主题设置 \tabularnewline
\hline
\end{tabularx}
\end{center}

\begin{center}\small\color{gray}SwiftUI 状态工具全景（2026 年，@Observable 已是主流）\end{center}

\begin{lstlisting}[language=Swift]
// iOS 17+ 的现代写法：@Observable 宏（推荐新项目使用）
import Observation

@Observable
class CounterModel {                 // 视图模型
    var count = 0
    func increment() { count += 1 }
}

struct CounterView: View {
    @State private var model = CounterModel()   // 持有模型

    var body: some View {
        VStack {
            Text("你点了 \(model.count) 次")
            Button("点我") { model.increment() }
        }
    }
}
\end{lstlisting}

给新手的粗暴建议：\textbf{页面私有的用 @State，跨页面/有逻辑的用 @Observable 类}。老教程里的 @StateObject/@ObservableObject 是上一代写法，能读懂即可，新代码按上面来。

\booksection{二、NavigationStack：页面跳转的“栈”}

iOS 16 起，多页面导航用\textbf{NavigationStack}——像一摞盘子，推一个进一个，返回时逐个弹出：

\begin{lstlisting}[language=Swift]
struct HomeView: View {
    var body: some View {
        NavigationStack {
            List {
                NavigationLink("张三的账本") {
                    DetailView(name: "张三的账本")   // 目标页面
                }
                NavigationLink("李四的账本") {
                    DetailView(name: "李四的账本")
                }
            }
            .navigationTitle("账本列表")   // 顶部大标题
        }
    }
}

struct DetailView: View {
    let name: String
    var body: some View {
        Text("\(name) 的详情页")
            .navigationTitle(name)
            .navigationBarTitleDisplayMode(.inline)
    }
}
\end{lstlisting}

NavigationLink 就是“能点的跳转链接”；\texttt{navigationTitle}自动生成导航栏标题和返回按钮。要\textbf{编程式跳转}（比如点完“保存”自动返回）时，用\texttt{@Environment(\textbackslash{}.\allowbreak dismiss)}关闭当前页面——这两招配合，90\% 的导航需求就齐了。

\booksection{三、弹窗与底部标签：sheet 与 TabView}

除了“推页面”，还有两种经典交互：

\begin{lstlisting}[language=Swift]
// 1. 底部弹出半屏/全屏（录入表单最爱）
struct HomeView: View {
    @State private var showAdd = false
    var body: some View {
        Button("记一笔") { showAdd = true }
            .sheet(isPresented: $showAdd) {
                AddRecordView()     // 弹出的录入页
            }
    }
}

// 2. 底部标签栏（App 的"骨架"）
struct MainTabView: View {
    var body: some View {
        TabView {
            RecordListView()
                .tabItem { Label("账本", systemImage: "book") }
            StatsView()
                .tabItem { Label("统计", systemImage: "chart.pie") }
            SettingsView()
                .tabItem { Label("设置", systemImage: "gear") }
        }
    }
}
\end{lstlisting}

\begin{tcbinfo}{小贴士：什么时候用哪种？}
经验法则：\textbf{主流程用 NavigationStack 推页面（列表详情），临时操作用 sheet 弹窗（新增、编辑、确认），并列的几大块用 TabView（底部标签）}。三者组合，就是绝大多数 App 的骨架。

\end{tcbinfo}

\booksection{四、List：iOS 的“国民列表”}

iOS 上最常用、最有“苹果味”的组件就是\textbf{List}——天然带分割线、滑动删除、点击高亮：

\begin{lstlisting}[language=Swift]
struct RecordListView: View {
    @State private var records = ["午饭 25.5", "地铁 6.0", "工资 +8000"]

    var body: some View {
        List {
            ForEach(records, id: \.self) { r in
                Text(r)
            }
            .onDelete { indexSet in          // 左滑删除
                records.remove(atOffsets: indexSet)
            }
        }
    }
}
\end{lstlisting}

列表 + NavigationLink + sheet + TabView——把这四样组合起来，你已经能搭出“账本列表 点开详情 点+弹录入”的完整 App 骨架了。下一章的记账本项目，就是这么搭的。

\begin{tcbpractice}{来动手：搭一个"三页骨架"}
\booksection{五、程序生命周期：前台与后台}

和安卓的 Activity 生命周期类似，iOS App 也有“出生、前台、后台、死亡”的状态变化。SwiftUI 里用\texttt{scenePhase}环境值观察：

\begin{lstlisting}[language=Swift]
@main
struct MyApp: App {
    @Environment(\.scenePhase) private var phase

    var body: some Scene {
        WindowGroup {
            ContentView()
        }
        .onChange(of: phase) { _, newPhase in
            switch newPhase {
            case .active:    print("回到前台，恢复刷新")
            case .inactive:  print("过渡中（来电/切后台瞬间）")
            case .background: print("进后台，保存草稿")
            @unknown default: break
            }
        }
    }
}
\end{lstlisting}

典型用途：\textbf{回前台时刷新数据，进后台时保存草稿}。iOS 对后台运行管得很严——普通 App 进后台后很快被挂起（冻结），别指望在后台偷偷干重活；真要定时任务，用第 68 章的本地通知，而不是“后台常驻”。

做一个假数据版本的记账本骨架：TabView 三个标签（账本/统计/设置） 账本页用 List 显示假数据、点行跳详情（NavigationStack） 右上角“+”按钮弹录入 sheet。不存数据、不联网，纯练骨架——这个骨架做好，第 65 章往里填真数据就是水到渠成。

\end{tcbpractice}

\begin{tcbsummary}{本章小结}
\begin{itemize}
  \item 状态四工具：@State（页面私有）、@Binding（借用）、@Observable（逻辑）、.environment（全局）
  \item 新项目用 @Observable 宏写视图模型；老教程的 @ObservableObject 能读即可
  \item NavigationStack 推页面；@Environment(\textbackslash{}.dismiss) 编程式返回
  \item 主流程推页面、临时操作弹 sheet、并列大块用 TabView
  \item List 是"国民列表"：分割线、滑动删除开箱即用
\end{itemize}

\end{tcbsummary}
